\section{An effective two-power conductor bound and original strict vanishing}
\label{sec:effective-two-power-strict}

\paragraph{The two coefficient fields play different roles.}
The finite field used in the Fourier argument below contains the values
of a \emph{fixed} character and fourth roots of unity.  It does not contain
all values of the varying two-power character.  Those values may belong
to an extension of arbitrarily large degree.  The original coefficient
ring of the strict group is introduced only after the uniform
nonvanishing argument.  Replacing the fixed field by that original
residue field in the degree argument would invalidate the bound.

\begin{theorem}[Effective two-power nonvanishing]
\label{G173-effective-two-power-nonvanishing}
Let $\ell$ be an odd prime and fix an embedding of the algebraic numbers into
$\overline{\mathbb Q}_\ell$.  Let $\theta$ be an odd finite Dirichlet
character of conductor $f_\theta$.  Suppose its reductions take values
in $k=\mathbb F_Q$, and that $\mu_4\subset k$.  Put
\[
 n_0=v_2(Q-1),\qquad L=\operatorname{lcm}(f_\theta,2^{n_0}).
 \tag{G173.1}
\]
For every even primitive Dirichlet character $\psi$ of conductor $2^c$,
\[
 \begin{gathered}
 c\geq2n_0,\qquad c>v_2(4f_\theta),\qquad
                       2^{c-n_0}>L-1\quad\Longrightarrow\\
                     L(0,\theta\psi)\text{ is an }\ell\text{-adic unit}.
 \end{gathered}
 \tag{G173.2}
\]
The assertion includes the required integrality of the special value.
It does not assume integrality of a Bernoulli expression with a
possibly $\ell$-divisible denominator.
\end{theorem}

\begin{proof}
\proofstep{1. The exact measure and special-value normalization.}
We give the complete specialization of Sinnott's argument needed here.
Its source is \emph{On a theorem of L.~Washington}, Ast\'erisque
147--148 (1987), printed pp.~215--217, Lemma~3.6 and its proof,
and pp.~220--222, Proposition~4.6 and its proof.
The tower prime in that paper is two here; its coefficient prime is
$\ell$.  Set $f=4f_\theta$ and
\[
 R(Z)=\frac{\displaystyle
       \sum_{\substack{1\leq a\leq f/2\\a\text{ odd}}}\theta(a)Z^a}
                                      {1-Z^f}.
 \tag{G173.3}
\]
In particular the denominator is $1-Z^f$, not $1-Z^{f/2}$.
For every two-power root of unity $\zeta$ in
$\overline{\mathbb F}_\ell$ prescribe
\[
 \widehat\alpha(\zeta)=
 \begin{cases}\overline R(\zeta),&\zeta^f\ne1,\\0,&\zeta^f=1.
 \end{cases}
 \tag{G173.4}
\]
Fourier inversion on each finite quotient of $\mathbb Z_2$ produces
compatible finitely additive measures: every power of two is
invertible in characteristic $\ell$.  The values lie in the fixed
field $k$.  Indeed (G173.4) is Frobenius equivariant over $k$, including
its zero exceptional values, so its finite inverse Fourier sums
are fixed by that Frobenius.
Since $R(-Z)=-R(Z)$ and the exceptional set is stable under
$\zeta\mapsto-\zeta$, the measure is supported on
$\mathbb Z_2^\times$.

For the character $\psi$ define
\[
 \tau(\psi,\zeta)=2^{-c}
       \sum_{a\bmod2^c}\psi(a)\zeta^{-a}.
 \tag{G173.5}
\]
The character is extended by zero on even integers.
Its primitive Gauss sums vanish at imprimitive $2^c$th roots.
Finite Fourier inversion gives
$\sum_{\zeta\text{ primitive}}\tau(\psi,\zeta)\zeta^a=\psi(a)$.
Oddness of $\theta$ gives the rational identity
\[
 R(Z)+R(Z^{-1})
  =\frac{\displaystyle\sum_{\substack{1\leq a\leq f\\a\text{ odd}}}
                                     \theta(a)Z^a}{1-Z^f}.
 \tag{G173.6}
\]
Evenness of $\psi$ gives
$\tau(\psi,\zeta^{-1})=\tau(\psi,\zeta)$.
Multiplying (G173.6) by the Gauss sums, summing over primitive roots,
and using the generating-series formula for $L(0,\theta\psi)$ yields
\[
 2\sum_{\zeta\text{ primitive}}\tau(\psi,\zeta)R(\zeta)
                                  =L(0,\theta\psi).
 \tag{G173.7}
\]
For clarity, the right side is the value at $Z=1$ of
$\sum_{a=1}^{f2^c}\theta(a)\psi(a)Z^a/(1-Z^{f2^c})$.
The product character is odd and nontrivial, so the singularity is
removable.  Its conductor is divisible by $2^c$, since
$c>v_2(f)$; consequently there is no missing Euler factor at two.
This proves precisely Sinnott's factor $1/2$.

Every denominator $1-\zeta^f$ in (G173.7) is an $\ell$-adic unit:
reduction preserves the order of a two-power root when $\ell\ne2$,
and $c>v_2(f)$.  The Gauss-sum denominator is a power of two.
Thus (G173.7) proves integrality, and after reduction gives
\[
 \Gamma_\alpha(\overline\psi)
 :=\int_{\mathbb Z_2^\times}\overline\psi(x)\,d\alpha(x)
             =\overline{\tfrac12L(0,\theta\psi)}.
 \tag{G173.8}
\]
The character values in this integral may lie in a larger finite
field than $k$; the measure itself is still $k$-valued.

\proofstep{2. A zero transform forces a zero at a root of full order.}
Write $U=1+4\mathbb Z_2$ and set $\beta=\alpha+\alpha\circ(-1)$,
where $(\alpha\circ(-1))(X)=\alpha(-X)$.
Then $\Gamma_\alpha(\overline\psi)=\int_U\overline\psi\,d\beta$.
Assume this is zero.  Put $n=c-n_0$, so $n\geq n_0\geq2$.
Choose $\zeta$ such that
\[
       \zeta^{2^n}=\overline\psi(1+2^n).
 \tag{G173.9}
\]
The right side has exact order $2^{n_0}$; hence $\zeta$ has
exact order $2^c$, not $2^{c-n_0}$.
The elementary two-adic lifting identity gives
\[
 [k(\mu_{2^c}):k]=2^{c-n_0},\qquad
 \operatorname{Gal}(k(\mu_{2^c})/k)
          =1+2^{n_0}\mathbb Z_2\pmod {2^c}.
 \tag{G173.10}
\]
Here the action is $\sigma_t(\zeta)=\zeta^t$.
For a root $\lambda\in\mu_{2^c}$, its normalized trace is
$\lambda$ if $\lambda\in\mu_{2^{n_0}}$, and zero otherwise;
this is the finite geometric sum over $t=1+2^{n_0}j$.

Multiply the vanishing integral by $\overline\psi(y)^{-1}$ for
$y\in U$ and take this normalized trace.  Because $\beta$ is
$k$-valued, it follows that
\[
 0=\int_{y(1+2^n\mathbb Z_2)}
                       \overline\psi(x/y)\,d\beta(x).
 \tag{G173.11}
\]
If $x/y=1+2^nz$, then
$(1+2^n)^z\equiv1+2^nz\pmod {2^c}$, since $2n\geq c$.
Thus $\overline\psi(x/y)=\zeta^{x/y-1}$, and (G173.11) also
holds with integrand $\zeta^{x/y}$.
Replace $y$ by $yt$ for $t\in1+2^{n_0}\mathbb Z_2$,
apply $\sigma_t$, and sum over representatives modulo
$1+2^n\mathbb Z_2$.  This proves
\[
 \int_{y(1+2^{n_0}\mathbb Z_2)}\zeta^{x/y}\,d\beta(x)=0.
 \tag{G173.12}
\]
This is the complete two-adic proof of the applied Lemma~3.6;
no qualitative ``infinitely many zeros'' conclusion is substituted
for its actual one-character assertion.

\proofstep{3. The restricted numerator is nonzero and has bounded degree.}
Take $y=1$ in (G173.12).  Restriction to
$1+2^{n_0}\mathbb Z_2$ is the finite root filter
\[
 \widehat{\beta|_{1+2^{n_0}\mathbb Z_2}}(\zeta)
 =2^{-n_0}\sum_{\epsilon^{2^{n_0}}=1}
                         \epsilon^{-1}\widehat\beta(\epsilon\zeta).
 \tag{G173.13}
\]
Every $\epsilon\zeta$ still has exact order $2^c$, since
$c\geq2n_0>n_0$.  In particular none belongs to the exceptional
set in (G173.4).  The equality with the rational function is exact
at every term of (G173.13), not merely valid at almost all roots.
By (G173.6), its formal series is
$\sum_{a\geq1,\ a\text{ odd}}\overline\theta(a)Z^a$.
Filtering that series gives exactly
\[
 \begin{gathered}
 \widehat{\beta|_{1+2^{n_0}\mathbb Z_2}}(\zeta)
              =\frac{P_1(\zeta)}{1-\zeta^L},\\
 P_1(Z)=\sum_{\substack{1\leq a\leq L\\a\equiv1\bmod2^{n_0}}}
                              \overline\theta(a)Z^a\in k[Z].
 \end{gathered}
 \tag{G173.14}
\]
The denominator is nonzero: $v_2(L)=\max(v_2(f_\theta),n_0)<c$.
The polynomial has degree at most $L-1$, and its coefficient of
$Z$ is $\overline\theta(1)=1$.  Reduction may lower the conductor
of $\theta$, but its period still divides $f_\theta$ and this
coefficient remains one.
By (G173.10), a primitive $2^c$th root has minimal-polynomial degree
$2^{c-n_0}$ over $k$.  The inequality in (G173.2) therefore makes
$P_1(\zeta)=0$ impossible.  Equations (G173.12)--(G173.14) contradict
the assumed vanishing of (G173.8), proving the theorem.
\end{proof}

\paragraph{The actual coefficient prime and finite labels.}
From now on put $p=269$ and $K=\mathbb Q(\sqrt{-3})$.
Let $\varepsilon$ be any of the six Dirichlet characters modulo seven,
including the trivial character with primitive conductor one.
For a primitive Dirichlet character $\delta$ of exact conductor $2^c$
define the sufficient threshold
\[
 M(\varepsilon)=
 \begin{cases}
 14,&\varepsilon=1,\\
 17,&\operatorname{ord}(\varepsilon)=2,\\
 19,&\operatorname{ord}(\varepsilon)=3\text{ or }6.
 \end{cases}
 \tag{G173.15}
\]
Use the same notation for their restrictions to $G_K$.
For a discrete coefficient $D$, let
\[
 \begin{gathered}
 S_{\rm str}(K,D)=\ker\bigl(H^1_{\rm cts}(G_K,D)
                                      \longrightarrow\mathcal L(D)\bigr),\\
 \mathcal L(D)=\prod_{v\mid p}H^1(G_{K_v},D)
                         \times\prod_{v\nmid p,\infty}H^1(I_v,D).
 \end{gathered}
 \tag{G173.16}
\]
In particular the local condition above $p$ is the full cohomology group.

\begin{theorem}[AC-G173: high two-power strict vanishing on original coefficients]
\label{AC-G173}
Assume $c\geq M(\varepsilon)$.
Let $E/\mathbb Q_p$ be any finite extension whose original ring of
integers $O$ contains the values of $\delta\varepsilon$; no bound on
the absolute ramification index is imposed.
For every $s\in\mathbb Z$, $e\geq1$, and continuous character
$\kappa:G_K\to1+\mathfrak m_O$ unramified at every finite place
away from $p$, one has
\[
 \begin{gathered}
 D=(E/O)(\delta\varepsilon\,\omega_p^{-s}\kappa^{-1}c_p^{-e}),\\
                             S_{\rm str}(K,D)=0.
 \end{gathered}
 \tag{G173.17}
\]
Here $\omega_p$ is the Teichmuller part of the cyclotomic character
$c_p$.  The residual strict group is also zero over every finite
coefficient field containing the residual character values.
There is no exceptional residue of $e+s$ modulo $268$.
For the actual host application, take $s=0$ and $\kappa=1$:
the finite label is exactly $\delta\varepsilon$, with trivial
cubic part.  No finite character at $269$ is added to the host.
\end{theorem}

\begin{proof}
\proofstep{1. Uniform odd special values over the fixed small field.}
Every primitive two-power character under consideration has a unique
factorization
\[
              \delta=\chi_{-4}^{\,a}\psi,
       \qquad a\in\{0,1\},\quad\psi(-1)=1,
 \tag{G173.18}
\]
where $\psi$ still has exact conductor $2^c$.
For $b=0,1$ and $j\bmod268$, the fixed characters needed below are
\[
              \theta=\chi_{-4}^{\,a}\chi_{-3}^{\,b}
                                      \varepsilon\omega_p^j.
 \tag{G173.19}
\]
Only the odd ones are used in the special-value theorem.
Their inverses and the fixed factors arising under reflection have
the same conductor bounds and the same order of $\varepsilon$.
The following table gives fixed fields containing these reductions;
it does not describe the growing original coefficient field.
\[
 \begin{array}{c|c|c|r|c}
 \varepsilon&k&n_0&\text{upper bound for }L&M(\varepsilon)\\ \hline
 1&\mathbb F_{269}&2&3228&14\\
 \operatorname{ord}=2&\mathbb F_{269}&2&22596&17\\
 \operatorname{ord}=3,6&\mathbb F_{269^2}&3&45192&19.
 \end{array}
 \tag{G173.20}
\]
Indeed $f_\theta$ divides $4\cdot3\cdot269$ in the first row and
$4\cdot3\cdot7\cdot269$ otherwise.
The cubic roots supplied by the order-three and order-six characters
modulo seven require $\mathbb F_{269^2}$, since $269\equiv-1\pmod3$.
Furthermore $v_2(269-1)=2$ and $v_2(269^2-1)=3$.
The respective minimal degrees at the thresholds are
$4096$, $32768$ and $65536$, larger than the displayed bounds.
All conditions in (G173.2), including $c>v_2(4f_\theta)$, hold.
Thus every odd character of the form (G173.19) times $\psi$ has
$p$-unit $L$-value at zero, as do the required inverse characters.
This statement holds for each embedding and survives any subsequent
extension of the coefficient residue field.

\proofstep{2. Integral Stickelberger operators kill the odd class sectors.}
Set
\[
 H=\mathbb Q(\zeta_{3\cdot7\cdot2^c p}),\qquad
 G=\operatorname{Gal}(H/\mathbb Q),\qquad
 J=\operatorname{Gal}(H/K),\qquad
 A_H=\operatorname{Cl}(H)\{p\}.
 \tag{G173.21}
\]
Its degree is $12\cdot2^{c-1}(p-1)$, prime to $p$.
Use a finite unramified coefficient extension containing all
characters of $G$ to form integral character idempotents.
Consider an odd character $\nu$ of $G$ whose two-primary conductor
is exactly $2^c$, whose conductor-seven component is
$\varepsilon$ or $\varepsilon^{-1}$, and whose remaining factors
are at three and $p$ as in (G173.19).
Let $f_\nu$ be its \emph{primitive} conductor and
$H_\nu=\mathbb Q(\zeta_{f_\nu})$.
The integral Stickelberger operator
\[
 (b_0-\sigma_{b_0})\Theta_{f_\nu},\qquad
 \Theta_f=f^{-1}\sum_{\substack{1\leq u\leq f\\(u,f)=1}}
                                      u\sigma_u^{-1},
 \tag{G173.22}
\]
annihilates $\operatorname{Cl}(H_\nu)$.
This is the exact integral convention of Sharifi,
Definition~4.5.1, Lemma~4.5.2 and Theorem~4.5.6;
no nonintegral isolated $\Theta_f$ is used.
Choose a positive $b_0$ by the Chinese remainder theorem with
\[
 b_0\equiv5\pmod {2^c},\qquad
 b_0\equiv1\pmod {3\cdot7\cdot p}.
 \tag{G173.23}
\]
The scalar of (G173.22) on the $\nu$ component is
\[
           (b_0-\nu(b_0))B_{1,\nu^{-1}}
              =-(b_0-\nu(b_0))L(0,\nu^{-1}).
 \tag{G173.24}
\]
The second factor is a unit by Step~1.  The reduction of the first
factor is $1-\overline\psi(5)$ or its inverse-character version.
It is nonzero: $5$ generates $1+4\mathbb Z/2^c\mathbb Z$ and
$\psi(5)$ has exact order $2^{c-2}$, which reduction preserves.
Consequently the entire $\nu$ class component of $H_\nu$ is zero.
Use of the primitive conductor here prevents an imprimitive
Euler factor from being mistaken for a unit.

This conclusion passes to the $\nu$ component of $H$.
For $V=\operatorname{Gal}(H/H_\nu)$, ideal extension $j$ and norm $N$
satisfy $Nj=[H:H_\nu]$ and $jN=\sum_{v\in V}v$ on class groups.
The degree is prime to $p$, and $\nu$ is trivial on $V$.
These identities identify the relevant class component with the
one at primitive conductor; in particular it is zero in $A_H$.
This is a component statement, not a claim that all of $A_H$ vanishes.

\proofstep{3. Reflection has no unramified root-of-unity correction.}
The exact specialized reflection argument retained in G168/G170
applies to this actual $H$; we spell out its field-dependent part.
The field is CM, abelian over $\mathbb Q$, of degree prime to $p$,
and contains $\mu_p$.
Let $B\subset H^\times$ be the radicals generating everywhere
unramified Kummer extensions.  The class-field/Kummer pairing and
$(x)=\mathfrak a^p$ give the equivariant sequence
\[
 0\longrightarrow
       (B\cap\mathcal O_H^\times)/\mathcal O_H^{\times p}
 \longrightarrow\operatorname{Hom}(A_H/pA_H,\mu_p)
 \longrightarrow A_H[p].
 \tag{G173.25}
\]
There is no surjectivity assertion at the last term.
The minus part of units modulo torsion is zero in a CM field;
its only possible odd unit line modulo $p$ is $\mu_p$, of character
$\omega_p$.  The $p$-primary roots in $H$ are precisely $\mu_p$.

At $p$, the conductor-$3\cdot7\cdot2^c$ factor is unramified.
Its residue degree is
\[
 \operatorname{lcm}(2,6,2^{c-2})=3\cdot2^{c-2},
 \qquad e(H_v/\mathbb Q_p)=p-1.
 \tag{G173.26}
\]
Here $269$ has order six modulo seven and exact order $2^{c-2}$
modulo $2^c$.  Adjoining $\mu_{p^2}$ has relative ramification
index $p$, unchanged by this unramified base change.
Thus the nonzero root-of-unity class cannot belong to the unramified
radical kernel in (G173.25).
Roots of unity of order prime to $p$ are already $p$th powers in $H$.
The odd kernel is zero, including its $\omega_p$ component.

For an even character $\nu$, the corresponding component of the
middle term in (G173.25) has character $\omega_p\nu^{-1}$.
It injects into $A_H[p](\omega_p\nu^{-1})$, and hence
\[
 A_H(\omega_p\nu^{-1})=0\quad\Longrightarrow\quad A_H(\nu)=0.
 \tag{G173.27}
\]
The reflected character is odd, still has exact two-primary
conductor $2^c$, and has the inverse conductor-seven component.
Its fixed factors still occur in (G173.19), so Step~2 kills it.
We have proved class zero in every required high-two-conductor
component, of either parity, without a complete class-number formula.

\proofstep{4. Strict residual classes select the original class character.}
Let $k'$ be any finite field containing the residual values, and put
\[
             \alpha=\overline\delta\,\overline\varepsilon
                              \overline\omega_p^{-r},
                     \qquad r\equiv e+s\pmod {268}.
 \tag{G173.28}
\]
Since $[H:K]$ is prime to $p$, restriction injects global $H^1$
into cohomology over $H$, where the coefficient is trivial.
For a cocycle $z$ and $g\in G_H$, $\tau\in G_K$, equivariance is
\[
                    z(\tau g\tau^{-1})=\alpha(\tau)z(g).
 \tag{G173.29}
\]
The strict local conditions make the resulting homomorphism vanish
on all inertia groups and on the full decomposition groups above $p$.
It factors through the class group with those primes removed, a
quotient of the ordinary class group.  Consequently
\[
 S_{\rm str}(K,k'(\alpha))\lhook\joinrel\longrightarrow
       \operatorname{Hom}_{\mathbb F_p[J]}
                            (A_H/pA_H,k'(\alpha)).
 \tag{G173.30}
\]
The character selected in (G173.30) is $\alpha$, not its inverse;
this is different from the duality in (G173.27).
Its two rational extensions differ by $\chi_{-3}$.
Both have the high two-primary conductor just treated, and
Steps~2--3 kill both class components.  If necessary, extend $k'$,
faithfully, to a finite splitting field for $G$; the integral
idempotents then make the right side zero, and this zero descends
back to $k'$.  This proves residual strict zero for all
$r$ and all $k'$, without identifying $k'$ with the fixed field in
(G173.20).

\proofstep{5. The original uniformizer and every local invariant.}
Let $\varpi$ be a uniformizer of the original $O$ and retain exactly
$D$ from (G173.17).  It is $\varpi$-divisible, and
$D[\varpi]\simeq k'(\alpha)$.
At $p$, $K_p/\mathbb Q_p$ is the unramified quadratic extension.
Since
\[
             v_2(269-1)=2,\qquad v_2(269+1)=1,
 \tag{G173.31}
\]
$269$ generates $1+4\mathbb Z/2^c\mathbb Z$.
Thus $\delta(269)$ has exact order $2^{c-2}$ and
$\delta(269)^2$ has exact order $2^{c-3}$.
In contrast $\varepsilon(269)^2$ has order dividing three.
It follows that
\[
      \overline\delta(\operatorname{Frob}_{K_p})
      \overline\varepsilon(\operatorname{Frob}_{K_p})
                 =\overline\delta(269)^2
                    \overline\varepsilon(269)^2\ne1.
 \tag{G173.32}
\]
Choose the Frobenius lift to fix $\mu_p$, which is possible because
$K_p$ is unramified and $\mathbb Q_p(\mu_p)$ is totally ramified.
Then the residual cyclotomic and Teichmuller factors are one on
this lift, for every $r$, and $\kappa$ reduces to one.
An actual scalar minus one is therefore an $O$-unit.
This proves $D^{G_{K_p}}=0$, including the residue $r=0$.
Using geometric instead of arithmetic Frobenius inverts the same
nontrivial scalar and does not change the argument.
In particular the global invariants are also zero.

At each $v\nmid p$, the cyclotomic, Teichmuller and $\kappa$ factors
are unramified.  The remaining inertia character has finite order
prime to $p$.  Its invariants in $D$ are either $D$ or zero:
a nontrivial root of unity of order prime to $p$ differs from one
by an $O$-unit.  In both cases these invariants are
$\varpi$-divisible.
Apply the cohomology sequence of
$0\to D[\varpi]\to D\xrightarrow{\varpi}D\to0$ globally,
on $G_{K_p}$, and on each away inertia group.
It follows that every strict element killed by $\varpi$ lifts to a
strict residual class.  The local maps from residual $H^1$ are
injective by the invariant calculation.  Step~4 gives
$S_{\rm str}(K,D)[\varpi]=0$.
Every continuous cocycle from a profinite group to the discrete
module $D$ has finite image, hence is annihilated by a power of
$\varpi$.  Therefore the whole strict group is zero.
This proof uses the uniformizer of the original ring; it imposes
no unramified-coefficient hypothesis and makes no change of lattice.
\end{proof}

\paragraph{Actual host characters and a single complete family count.}
In the prescribed hosts of (FG.3), the conductor-two cyclotomic
factor is $\mathbb Q(\zeta_{2^{h+2}})$ and the conductor-seven
factor is $\mathbb Q(\zeta_7)$, with $h\geq2$.
Every label $\delta\varepsilon$ above with $c\leq h+2$ therefore
occurs, with trivial cubic and trivial auxiliary quadratic factors.
They remain distinct on $G_K$: a conductor supported at two and
seven cannot become the nontrivial conductor-three character.
Their orders divide $3\cdot2^h$, so their values lie in the
original cyclotomic coefficient ring used in (FC.2).
No new host or finite $p$-character is introduced.

For integers $M\geq3$ define
\[
 B(h;M)=\max\{0,\,2^{h+1}-2^{M-2}\},\qquad
 C(h)=B(h;14)+B(h;17)+4B(h;19).
 \tag{G173.33}
\]
There are $2^{c-2}$ primitive Dirichlet characters of conductor
$2^c$, including both parities.  Summing their counts proves that
$C(h)$ is exactly the number of labels established by this theorem.
For $h\geq17$ it is
\[
                       C(h)=12\cdot2^h-561152.
 \tag{G173.34}
\]
For every integer cutoff $T\geq1$, all $C(h)(T-1)$ pairs consisting
of one such label and a weight $1\leq e<T$ have strict group zero
at $p=269$ on their original coefficient rings.
These high-conductor labels are disjoint from the earlier
conductor-at-most-sixteen label families.

The precise (FC.2) comparison is inflation from $G_{K,\Sigma}$
into $G_K$, with $s=0$ and $\kappa=1$ in (G173.17).
The coefficient is unramified outside the prescribed set.
Outside that set the inflated cocycle already kills inertia;
inside it the away conditions are the same, and the full local
condition at $269$ is unchanged.  Thus a (FC.2) strict class would
be a class in (G173.16), and is zero.
This is a literal original-coefficient statement; the separate
positive-index/local correction bound and common-window assembly
remain those of (FL.4), (EU.2) and the actual window theorem.
Strict zero alone is not asserted to make a derived index zero.

\paragraph{Scope of the uniform improvement.}
The conductor thresholds are independent of the varying auxiliary
primes $q,d$, of the prescribed window, and of every weight below
the cutoff $T_{399}$.  The proof supplies actual strict zero on any
such window; it does not need to choose a new one.
It treats the specified trivial-cubic labels only.
Multiplication by the growing auxiliary quadratic characters
would enlarge $f_\theta$ by their conductors; (G173.2) would then
require a correspondingly larger two-power level.  In the actual
family their logarithmic conductors can exceed the available
$h=O(\log u)$ scale, so this argument does not cover those labels.
The abstract nonvanishing theorem is valid at every odd coefficient
prime, but no all-prime strict statement is inferred: the
prime-to-coefficient group degrees, local invariants and original
integral comparisons would require their own checks.
No mixed-cubic estimate, exceptional-weight theorem for mixed terms,
higher-regulator bound, or complete $I_{399}$ budget follows here.

\paragraph{Geometric dictionary and source scope.}
The finite abelian characters correspond to rank-one local systems
on the finite cyclotomic covers after their ramified primes are
removed.  The unramified Kummer radicals in (G173.25) are the
$\mu_p$-torsors of the retained Kummer/class-field dictionary;
their unit kernel is checked at every prime above $p$.
This uses the already audited SGA~1, Exp.~V, \S\S4,7 and
Stacks~03PK scopes, not a new scan attribution.
The new quantitative argument uses the complete primary
Lemma~3.6 and Proposition~4.6 proofs cited above; the applied
Stickelberger and reflection inputs retain the full proof scopes
recorded for G168/G170.  The uniform strict result has no endpoint,
automorphy, native support, model-ring or regulator premise.
